\documentclass[11pt,hacksoc]{meetingmins}
\setcommittee{HackSoc <> School of Computer Science}

\setdate{June 17, 2019}

\setpresent{
  \chair{Max~Willson},
  Christian~Wagner,
  Mike~Pound,
  Tim~Miller,
  Hani~Momenmia,
  Aaron~Osher,
  Primoz~Fabiani
}

\begin{document}
\maketitle

\section{HackSoc \& Cyber Security in 2019/20}
\begin{items}
	\item Activities that Hani Would Like to do in 2019/20
		\subitem More Workshops
		\subitem Capture the Flag Competitions
		\subitem Cyber Team\footnote{The Cyber sub-group shall be a group of members with an exceptional interest in cybersecurity and penetration testing, led by the Cyber Security Secretary, that are committed to competing in CTF competitions, as well as attending conferences and meet-ups.} (referred to as the Cyber Sub Group in the Constitution)
		
	\item Capture The Flag Competition
		\subitem Hani recommends running everything in Docker containers to ensure the sandbox environment is easily isolated
		\subitem Wants to use Hack the Box or similar
			\subsubitem Won't work because: The network blocks the ports required to use external services
			\subsubitem They have dedicated servers that Universities can use\footnote{https://www.hackthebox.eu/universities}
			\subsubitem To get access to the platform, there is paper work that the University needs to fill out
			\subsubitem Hack the Box has access to servers in browser
			\subsubitem It's reasonable for HackSoc to run a Hack the Box CTF event, a proposal (date, resources needed etc.) would need to be shared with the School and Information Services
		\subitem Spoke with Mark Goodwin at Mozilla about CTFs that used to run
			\subsubitem Mozilla said they could sponsor any CTFs
			\subsubitem Also spoke with Microsoft about this
		\subitem Wants to create a platform that allows challenges to be created and run in Docker Containers
			\subsubitem Will take approx. 1 year to develop
			\subsubitem Could be created as a Level 2 Group Project 
			\subsubitem Should it become a group project, a sponsor (perhaps Mozilla) would have to exist
			\subsubitem Said project needs to have some sort of a written proposal/specification
			\subsubitem Is working with Derby Uni and NTU Students
			
	\item Workshops
		\subitem Wants workshops to compliment the skills needed to participate in the CTFs
		
	\item Cyber Team
		\subitem No issue raised with respect to taking people to events - SU is the main PoC
		\subitem Don't mention "Penetration Testing" to Information Services
		\subitem Unlikely to have any issues with doing Penetration Testing
		
	\item Christian is working with the NCSC\footnote{National Cyber Security Centre} on vulnerability testing
		\subitem Can potentially reach out to help support events and strengthen training that we can do
		
	\item Primary contacts when getting stuff approved
		\subitem Information Services (IS)
			\subsubitem School can potentially facilitate communicating with them
			\subsubitem Likely to not always be the most friendly with respect to working with HackSoc, especially considering they're constantly stressed with ensuring up time of systems
		\subitem Research Systems Team (RST)\footnote{School of Computer Science's Dedicated IT Team}
		\subitem We should create a ethics agreement 
		
	\item School's Cyber Security Lab/Activites
		\subitem Having an isolated subnet or dedicated network
		\subitem HackSoc to share what we'd like to have in a Cyber Security Lab
		\subitem Could potentially expand on top of the labs done in G53SEC
			\subsubitem Mike will provide a list of labs done in G53SEC
		
	\item Future Liaising between HackSoc and the School
		\subitem Future proposals for events would need to be passed via Max, and then would be communicated to the appropriate individuals
		
\end{items}

\end{document}